\documentclass[12pt,a4paper]{article}

% ── Packages ──
\usepackage[utf8]{inputenc}
\usepackage[T1]{fontenc}
\usepackage{amsmath,amssymb}
\usepackage{geometry}
\usepackage{graphicx}
\usepackage{booktabs}
\usepackage{array}
\usepackage{enumitem}
\usepackage{hyperref}
\usepackage{xcolor}
\usepackage{listings}
\usepackage{fancyhdr}
\usepackage{titlesec}
\usepackage{tcolorbox}

\geometry{margin=2.5cm}
\hypersetup{colorlinks=true, linkcolor=blue!60!black, urlcolor=blue!60!black}

% ── Header / Footer ──
\pagestyle{fancy}
\fancyhf{}
\fancyhead[L]{\textit{Rotations}}
\fancyhead[R]{\textit{GNC Fundamentals}}
\fancyfoot[C]{\thepage}

% ── Custom boxes ──
\newtcolorbox{keybox}[1][]{colback=blue!5!white, colframe=blue!50!black,
  fonttitle=\bfseries, title=#1}

\newtcolorbox{warningbox}[1][]{colback=red!5!white, colframe=red!50!black,
  fonttitle=\bfseries, title=#1}

% ── MATLAB listing style ──
\lstdefinestyle{matlab}{
  language=Matlab,
  basicstyle=\ttfamily\small,
  keywordstyle=\color{blue},
  commentstyle=\color{green!50!black},
  stringstyle=\color{red!70!black},
  numbers=left,
  numberstyle=\tiny\color{gray},
  frame=single,
  breaklines=true,
  captionpos=b
}

% ── Title ──
\title{\textbf{Rotations: Euler Angles, Gimbal Lock,\\[0.3em]
       and Quaternions}\\[0.5em]
       \large A Comprehensive Guide for GNC Engineers}
\author{}
\date{April 2026}

\renewcommand{\contentsname}{Index}

\begin{document}
\maketitle
\tableofcontents
\newpage

% ══════════════════════════════════════════════════════════════════════════════
\section{Euler Angle Rotations}
% ══════════════════════════════════════════════════════════════════════════════

Any orientation in 3D space can be described by a sequence of three successive
rotations about coordinate axes. These are called \textbf{Euler angles}. The
two fundamental conventions for defining these rotations are
\emph{intrinsic} and \emph{extrinsic}.

% ──────────────────────────────────────────────────────────────────────────────
\subsection{Rotation Convention: Right-Hand Rule}
% ──────────────────────────────────────────────────────────────────────────────

This document uses the \textbf{right-hand rule} throughout, which is the
standard convention in aerospace, GNC, and most engineering disciplines.

\begin{keybox}[Right-Hand Rule]
Point your \textbf{right thumb} along the positive direction of the rotation
axis. Your fingers curl in the direction of \textbf{positive rotation}
(counter-clockwise when viewed from the tip of the axis).
\end{keybox}

\begin{table}[h]
\centering
\begin{tabular}{@{}lll@{}}
\toprule
 & \textbf{Right-Hand Rule} & \textbf{Left-Hand Rule} \\
\midrule
Positive rotation & Counter-clockwise (looking down axis) &
  Clockwise (looking down axis) \\
Cross product     & $\hat{x}\times\hat{y}=+\hat{z}$ &
  $\hat{x}\times\hat{y}=-\hat{z}$ \\
Used in           & Aerospace, physics, most GNC &
  Some computer graphics (e.g.\ DirectX) \\
Matrix form       & $\sin\theta$ terms as shown below &
  $\sin\theta$ terms negated \\
\bottomrule
\end{tabular}
\caption{Right-hand vs.\ left-hand rotation conventions.}
\end{table}

\begin{warningbox}[Convention matters]
Mixing right-hand and left-hand conventions silently negates rotation
directions. Always verify which convention your codebase, sensor, and
simulation tool use before combining rotations.
\end{warningbox}

% ──────────────────────────────────────────────────────────────────────────────
\subsection{Elementary Rotation Matrices}
% ──────────────────────────────────────────────────────────────────────────────

Using the right-hand rule, the three elementary (principal-axis) rotation
matrices are:

\subsubsection*{Rotation about the X-axis by angle $\phi$ (roll)}

\begin{equation}
  R_X(\phi) = \begin{pmatrix}
    1 & 0          & 0           \\
    0 & \cos\phi  & \sin\phi   \\
    0 & -\sin\phi & \cos\phi
  \end{pmatrix}
  \label{eq:Rx}
\end{equation}

\subsubsection*{Rotation about the Y-axis by angle $\theta$ (pitch)}

\begin{equation}
  R_Y(\theta) = \begin{pmatrix}
    \cos\theta  & 0 & -\sin\theta \\
    0            & 1 & 0            \\
    \sin\theta  & 0 & \cos\theta
  \end{pmatrix}
  \label{eq:Ry}
\end{equation}

\subsubsection*{Rotation about the Z-axis by angle $\psi$ (yaw)}

\begin{equation}
  R_Z(\psi) = \begin{pmatrix}
    \cos\psi  & \sin\psi  & 0 \\
    -\sin\psi & \cos\psi  & 0 \\
    0          & 0           & 1
  \end{pmatrix}
  \label{eq:Rz}
\end{equation}

\begin{keybox}[Sign pattern]
In the right-hand convention for nav-to-body matrices, the $\sin$ term with
the \textbf{negative sign} appears:
\begin{itemize}
  \item $R_X$: row 3, column 2 \quad ($-\sin\phi$)
  \item $R_Y$: row 1, column 3 \quad ($-\sin\theta$)
  \item $R_Z$: row 2, column 1 \quad ($-\sin\psi$)
\end{itemize}
For a left-hand convention, these signs flip.
\end{keybox}

% ──────────────────────────────────────────────────────────────────────────────
\subsection{Intrinsic Rotations}
% ──────────────────────────────────────────────────────────────────────────────

Intrinsic rotations are performed about the \textbf{axes of the rotating body}
(body-fixed frame). Each successive rotation occurs about an axis that has been
updated by the previous rotation.

For the \textbf{ZYX intrinsic sequence} (common in aerospace as
yaw--pitch--roll):

\begin{equation}
  R = R_Z(\psi)\; R_Y(\theta)\; R_X(\phi)
\end{equation}

\begin{enumerate}
  \item Rotate about the body's \textbf{Z} axis by $\psi$ (yaw).
  \item Rotate about the \textbf{new $Y'$} axis by $\theta$ (pitch).
  \item Rotate about the \textbf{new $X''$} axis by $\phi$ (roll).
\end{enumerate}

% ──────────────────────────────────────────────────────────────────────────────
\subsection{Extrinsic Rotations}
% ──────────────────────────────────────────────────────────────────────────────

Extrinsic rotations are performed about the \textbf{axes of the fixed/world
frame}. The axes never change between rotations.

For the \textbf{XYZ extrinsic sequence}:

\begin{equation}
  R = R_Z(\gamma)\; R_Y(\beta)\; R_X(\alpha)
\end{equation}

\begin{enumerate}
  \item Rotate about the fixed \textbf{X} axis by $\alpha$.
  \item Rotate about the fixed \textbf{Y} axis by $\beta$.
  \item Rotate about the fixed \textbf{Z} axis by $\gamma$.
\end{enumerate}

% ──────────────────────────────────────────────────────────────────────────────
\subsection{Key Equivalence}
% ──────────────────────────────────────────────────────────────────────────────

\begin{keybox}[Intrinsic--Extrinsic Equivalence]
An \textbf{intrinsic} rotation sequence in order $A$-$B$-$C$ produces the
\textbf{same} result as an \textbf{extrinsic} rotation in the
\textbf{reversed} order $C$-$B$-$A$ with the same angles:
%
\begin{equation}
  R_{\text{intrinsic}}(Z\text{-}Y\text{-}X)
  \;=\; R_{\text{extrinsic}}(X\text{-}Y\text{-}Z)
\end{equation}
\end{keybox}

This holds because:
\begin{itemize}
  \item \textbf{Intrinsic}: multiply rotation matrices left-to-right in the
        order applied $\;\Rightarrow\; R_Z \cdot R_Y \cdot R_X$.
  \item \textbf{Extrinsic}: multiply rotation matrices right-to-left in the
        order applied $\;\Rightarrow\;$ same product $R_Z \cdot R_Y \cdot R_X$.
\end{itemize}

\begin{table}[h]
\centering
\begin{tabular}{@{}lll@{}}
\toprule
 & \textbf{Intrinsic} & \textbf{Extrinsic} \\
\midrule
Axes          & Body-fixed (rotate with body)  & World-fixed (stationary) \\
Matrix order  & Multiply in sequence order     & Multiply in reverse order \\
Common use    & Aerospace (yaw--pitch--roll)   & Animation, some robotics \\
Gimbal lock   & Middle rotation $= \pm 90°$   & Same condition, same angles \\
\bottomrule
\end{tabular}
\caption{Comparison of intrinsic and extrinsic conventions.}
\end{table}

% ══════════════════════════════════════════════════════════════════════════════
\section{Nav-to-Body or Body-to-Nav?}
% ══════════════════════════════════════════════════════════════════════════════

With the standard intrinsic ZYX convention the composed matrix

\begin{equation}
  R_{\text{intrinsic}} = R_Z(\psi)\; R_Y(\theta)\; R_X(\phi)
\end{equation}

is the \textbf{nav-to-body} (world-to-body) rotation matrix. It transforms a
vector from the navigation/world frame into the body frame:

\begin{equation}
  \vec{v}_{\text{body}} = R_{\text{intrinsic}}\;\vec{v}_{\text{nav}}
\end{equation}

To go \textbf{body-to-nav}, use the transpose (rotation matrices are
orthogonal, so $R^{-1} = R^T$):

\begin{equation}
  R_{\text{body}\to\text{nav}} = R_{\text{intrinsic}}^T
    = R_X(-\phi)\; R_Y(-\theta)\; R_Z(-\psi)
\end{equation}

% ──────────────────────────────────────────────────────────────────────────────
\subsection{Expanded Nav-to-Body Matrix (ZYX)}
% ──────────────────────────────────────────────────────────────────────────────

Multiplying $R_Z(\psi)\,R_Y(\theta)\,R_X(\phi)$ using the elementary
matrices from Equations~\eqref{eq:Rz},~\eqref{eq:Ry},~\eqref{eq:Rx}:

\begin{equation}
\boxed{
  R_{\text{nav}\to\text{body}} =
  \begin{pmatrix}
    \cos\theta\cos\psi &
    \cos\theta\sin\psi &
    -\sin\theta \\
    \sin\phi\sin\theta\cos\psi - \cos\phi\sin\psi &
    \sin\phi\sin\theta\sin\psi + \cos\phi\cos\psi &
    \sin\phi\cos\theta \\
    \cos\phi\sin\theta\cos\psi + \sin\phi\sin\psi &
    \cos\phi\sin\theta\sin\psi - \sin\phi\cos\psi &
    \cos\phi\cos\theta
  \end{pmatrix}
}
\label{eq:nav2body}
\end{equation}

where $\psi$ = yaw, $\theta$ = pitch, $\phi$ = roll (all right-hand
positive).

% ──────────────────────────────────────────────────────────────────────────────
\subsection{Expanded Body-to-Nav Matrix (ZYX)}
% ──────────────────────────────────────────────────────────────────────────────

Since $R$ is orthogonal, $R_{\text{body}\to\text{nav}} =
R_{\text{nav}\to\text{body}}^T$:

\begin{equation}
\boxed{
  R_{\text{body}\to\text{nav}} =
  \begin{pmatrix}
    \cos\theta\cos\psi &
    \sin\phi\sin\theta\cos\psi - \cos\phi\sin\psi &
    \cos\phi\sin\theta\cos\psi + \sin\phi\sin\psi \\
    \cos\theta\sin\psi &
    \sin\phi\sin\theta\sin\psi + \cos\phi\cos\psi &
    \cos\phi\sin\theta\sin\psi - \sin\phi\cos\psi \\
    -\sin\theta &
    \sin\phi\cos\theta &
    \cos\phi\cos\theta
  \end{pmatrix}
}
\label{eq:body2nav}
\end{equation}

\begin{keybox}[Quick check]
At zero attitude ($\psi=\theta=\phi=0$), both matrices reduce to the
$3\times3$ identity, confirming body frame $=$ nav frame.
\end{keybox}

\begin{keybox}[Why nav-to-body?]
The intrinsic ZYX sequence describes how to rotate \emph{from} the nav frame
\emph{into} the body frame: first yaw about nav Z, then pitch about the
intermediate $Y'$, then roll about the final body $X''$. The composed matrix
maps nav-frame coordinates into the resulting body frame.
\end{keybox}

% ──────────────────────────────────────────────────────────────────────────────
\subsection{Passive vs.\ Active --- Two Conventions, Same Product Order}
% ──────────────────────────────────────────────────────────────────────────────

The product $R_Z(\psi)\,R_Y(\theta)\,R_X(\phi)$ can yield \emph{either}
nav-to-body or body-to-nav, depending on how the elementary matrices are
defined. Both use the right-hand rule and the \textbf{same multiplication
order}; the only difference is the \textbf{sign of $\sin$} in each
elementary matrix.

% ─ ─ ─ ─ ─ ─ ─ ─ ─ ─ ─ ─ ─ ─ ─ ─ ─ ─ ─ ─ ─ ─ ─ ─ ─ ─ ─ ─ ─ ─ ─ ─ ─ ─ ─
\subsubsection*{GNC Aerospace Convention --- Passive (Alias)}
% ─ ─ ─ ─ ─ ─ ─ ─ ─ ─ ─ ─ ─ ─ ─ ─ ─ ─ ─ ─ ─ ─ ─ ─ ─ ─ ─ ─ ─ ─ ─ ─ ─ ─ ─

Standard aerospace textbooks (Titterton \& Weston, Stevens \& Lewis, etc.)
use the \textbf{passive} (alias) elementary matrices listed in
Equations~\eqref{eq:Rz}--\eqref{eq:Rx}:

\begin{equation}
  R_Z^{(P)}(\psi) = \begin{pmatrix}
    \cos\psi  & \sin\psi  & 0 \\
    -\sin\psi & \cos\psi  & 0 \\
    0          & 0          & 1
  \end{pmatrix}
  \quad (+\sin\text{ above diagonal})
  \label{eq:RzP}
\end{equation}

The composed product gives the \textbf{nav-to-body} direction-cosine matrix:

\begin{equation}
\boxed{
  R_Z^{(P)}\, R_Y^{(P)}\, R_X^{(P)}
  \;=\; C_n^b
  \quad\text{(nav-to-body)}
}
\end{equation}

Meaning: the matrix \emph{re-expresses} a stationary vector in the rotated
body frame.  The frame moves; the vector stays.

\[
  \vec{v}_{\text{body}} = C_n^b \;\vec{v}_{\text{nav}}
\]

% ─ ─ ─ ─ ─ ─ ─ ─ ─ ─ ─ ─ ─ ─ ─ ─ ─ ─ ─ ─ ─ ─ ─ ─ ─ ─ ─ ─ ─ ─ ─ ─ ─ ─ ─
\subsubsection*{CTCT Convention --- Active (Alibi)}
% ─ ─ ─ ─ ─ ─ ─ ─ ─ ─ ─ ─ ─ ─ ─ ─ ─ ─ ─ ─ ─ ─ ─ ─ ─ ─ ─ ─ ─ ─ ─ ─ ─ ─ ─

The CTCT codebase (e.g.\ \texttt{Rotation\_Matrix.m},
\texttt{DozerMinMaxYawCalTestHelpers.h}, autocode \texttt{Rot\_Mtrx\_YRP})
uses \textbf{active} (alibi) elementary matrices with the opposite $\sin$
sign.

\begin{warningbox}[CTCT Angle Naming]
CTCT uses non-standard angle names relative to aerospace:
\begin{center}
\begin{tabular}{@{}llll@{}}
\toprule
\textbf{Axis} & \textbf{Aerospace name} & \textbf{CTCT name} & \textbf{Symbol} \\
\midrule
Z & Yaw   & Yaw   & $\psi$ \\
Y & Pitch & Roll  & $\phi_R$ \\
X & Roll  & Pitch & $\theta_P$ \\
\bottomrule
\end{tabular}
\end{center}
The intrinsic sequence is therefore
\textbf{Yaw--Roll--Pitch} in CTCT terminology
vs.\ \textbf{Yaw--Pitch--Roll} in aerospace.
The matrix product $R_Z\,R_Y\,R_X$ is identical---only the labels differ.
\end{warningbox}

The three CTCT active elementary matrices are:

\subsubsection*{CTCT: Rotation about Z by $\psi$ (Yaw)}

\begin{equation}
  R_Z^{(A)}(\psi) = \begin{pmatrix}
    \cos\psi  & -\sin\psi & 0 \\
    \sin\psi  & \cos\psi  & 0 \\
    0          & 0          & 1
  \end{pmatrix}
  \quad (+\sin\text{ below diagonal})
  \label{eq:RzA}
\end{equation}

\subsubsection*{CTCT: Rotation about Y by $\phi_R$ (Roll)}

\begin{equation}
  R_Y^{(A)}(\phi_R) = \begin{pmatrix}
    \cos\phi_R  & 0 & \sin\phi_R \\
    0            & 1 & 0            \\
    -\sin\phi_R & 0 & \cos\phi_R
  \end{pmatrix}
  \label{eq:RyA}
\end{equation}

\subsubsection*{CTCT: Rotation about X by $\theta_P$ (Pitch)}

\begin{equation}
  R_X^{(A)}(\theta_P) = \begin{pmatrix}
    1 & 0              & 0               \\
    0 & \cos\theta_P  & -\sin\theta_P  \\
    0 & \sin\theta_P  & \cos\theta_P
  \end{pmatrix}
  \label{eq:RxA}
\end{equation}

\begin{keybox}[Active sign pattern]
Compare with the passive (aerospace) elementary matrices in
Equations~\eqref{eq:Rx}--\eqref{eq:Rz}: each active matrix is the
\textbf{transpose} of its passive counterpart.
The $-\sin$ terms move to the opposite off-diagonal position:
\begin{itemize}
  \item $R_X^{(A)}$: $-\sin\theta_P$ at row 2, col 3 \quad
        (passive: row 3, col 2)
  \item $R_Y^{(A)}$: $-\sin\phi_R$ at row 3, col 1 \quad
        (passive: row 1, col 3)
  \item $R_Z^{(A)}$: $-\sin\psi$ at row 1, col 2 \quad
        (passive: row 2, col 1)
\end{itemize}
\end{keybox}

Using the \textbf{same multiplication order} $R_Z\,R_Y\,R_X$ gives the
\textbf{body-to-nav} matrix:

\begin{equation}
\boxed{
  R_Z^{(A)}(\psi)\, R_Y^{(A)}(\phi_R)\, R_X^{(A)}(\theta_P)
  \;=\; C_b^n
  \quad\text{(body-to-nav)}
}
\end{equation}

Meaning: the matrix physically \emph{moves} a body-frame vector into the nav
frame.  The vector moves; the frame stays.

\[
  \vec{v}_{\text{nav}} = C_b^n \;\vec{v}_{\text{body}}
\]

\subsubsection*{Expanded CTCT Body-to-Nav Matrix}

Multiplying $R_Z^{(A)}(\psi)\,R_Y^{(A)}(\phi_R)\,R_X^{(A)}(\theta_P)$:

\begin{equation}
\boxed{
  C_b^n =
  \begin{pmatrix}
    \cos\phi_R\cos\psi &
    -\cos\theta_P\sin\psi + \sin\theta_P\sin\phi_R\cos\psi &
    \sin\theta_P\sin\psi + \cos\theta_P\sin\phi_R\cos\psi \\[4pt]
    \cos\phi_R\sin\psi &
    \cos\theta_P\cos\psi + \sin\theta_P\sin\phi_R\sin\psi &
    -\sin\theta_P\cos\psi + \cos\theta_P\sin\phi_R\sin\psi \\[4pt]
    -\sin\phi_R &
    \sin\theta_P\cos\phi_R &
    \cos\theta_P\cos\phi_R
  \end{pmatrix}
}
\label{eq:ctct-body2nav}
\end{equation}

where $\psi$ = Yaw, $\phi_R$ = Roll (CTCT, about Y), $\theta_P$ = Pitch
(CTCT, about X).  This is the \textbf{transpose} of the aerospace
nav-to-body matrix (Eq.~\ref{eq:nav2body}) with the CTCT naming applied.

% ─ ─ ─ ─ ─ ─ ─ ─ ─ ─ ─ ─ ─ ─ ─ ─ ─ ─ ─ ─ ─ ─ ─ ─ ─ ─ ─ ─ ─ ─ ─ ─ ─ ─ ─
\subsubsection*{Why the Same Order Gives Transposes}
% ─ ─ ─ ─ ─ ─ ─ ─ ─ ─ ─ ─ ─ ─ ─ ─ ─ ─ ─ ─ ─ ─ ─ ─ ─ ─ ─ ─ ─ ─ ─ ─ ─ ─ ─

Each active elementary matrix is the transpose of its passive counterpart:
$R_{\text{axis}}^{(A)} = \bigl(R_{\text{axis}}^{(P)}\bigr)^T$.

A composed product of individually-transposed factors is \textbf{not} the
same as transposing the whole product---the algebraic transpose rule
$(ABC)^T = C^T B^T A^T$ reverses the order.  So how can the same-order
active product equal the transpose of the same-order passive product?

The answer is that the expanded $3\times3$ matrix happens to satisfy the
identity:

\begin{equation}
\boxed{
  R_Z^{(A)}\, R_Y^{(A)}\, R_X^{(A)}
  \;=\;
  \bigl(R_Z^{(P)}\, R_Y^{(P)}\, R_X^{(P)}\bigr)^T
}
\label{eq:same-order-transpose}
\end{equation}

This can be verified element-by-element: flipping the sign of every $\sin$
term in every elementary matrix and then multiplying in the same order
produces exactly the transpose of the passive product.

\begin{keybox}[Numerical verification]
At $\psi=30^\circ,\;\theta=20^\circ,\;\phi=10^\circ$:
\begin{tabular}{@{}ll@{}}
$C_n^b = R_Z^{(P)} R_Y^{(P)} R_X^{(P)}$: &
  $(1,2) = +0.4699,\; (2,1) = -0.4410$ \\
$C_b^n = R_Z^{(A)} R_Y^{(A)} R_X^{(A)}$: &
  $(1,2) = -0.4410,\; (2,1) = +0.4699$ \\
\end{tabular}\\[4pt]
These are transposes of each other.  The reversed-order product
$R_X^{(A)} R_Y^{(A)} R_Z^{(A)}$ gives $(1,2) = -0.4699,\; (2,1) = +0.5438$
---a \textbf{completely different} matrix.
\end{keybox}

\subsubsection*{Reconciling with the Transpose Reversal Rule}

The algebraic identity $(ABC)^T = C^T B^T A^T$ is always true.  So:

\begin{equation}
  \bigl(R_Z^{(P)} R_Y^{(P)} R_X^{(P)}\bigr)^T
  = R_X^{(A)}\, R_Y^{(A)}\, R_Z^{(A)}
  \qquad\text{(reversed order, purely algebraic)}
\end{equation}

This reversed product $R_X^{(A)} R_Y^{(A)} R_Z^{(A)}$ is an $X$-$Y$-$Z$
rotation sequence---a \emph{different} set of physical rotations than
$Z$-$Y$-$X$.  Meanwhile, the same-order product
$R_Z^{(A)} R_Y^{(A)} R_X^{(A)}$ is a $Z$-$Y$-$X$ sequence with
individually-transposed factors.  Both happen to produce the \emph{same}
numerical matrix (which is $C_b^n$), but through different mathematical
paths.  There is no contradiction.

% ─ ─ ─ ─ ─ ─ ─ ─ ─ ─ ─ ─ ─ ─ ─ ─ ─ ─ ─ ─ ─ ─ ─ ─ ─ ─ ─ ─ ─ ─ ─ ─ ─ ─ ─
\subsubsection*{Side-by-Side Comparison}
% ─ ─ ─ ─ ─ ─ ─ ─ ─ ─ ─ ─ ─ ─ ─ ─ ─ ─ ─ ─ ─ ─ ─ ─ ─ ─ ─ ─ ─ ─ ─ ─ ─ ─ ─

\begin{table}[h]
\centering
\begin{tabular}{@{}p{3.8cm}p{5cm}p{5cm}@{}}
\toprule
 & \textbf{GNC Aerospace (Passive)} & \textbf{CTCT (Active)} \\
\midrule
Elementary $R_Z$ &
  $+\sin\psi$ above diagonal &
  $+\sin\psi$ below diagonal \\[4pt]
Product order &
  $R_Z\, R_Y\, R_X$ &
  $R_Z\, R_Y\, R_X$ \textbf{(same)} \\[4pt]
Result of product &
  $C_n^b$ (nav-to-body) &
  $C_b^n$ (body-to-nav) \\[4pt]
To get the other direction &
  Transpose: $C_b^n = (C_n^b)^T$ &
  Transpose: $C_n^b = (C_b^n)^T$ \\[4pt]
Key relationship &
  \multicolumn{2}{c}{$R_Z^{(A)} R_Y^{(A)} R_X^{(A)}
  = \bigl(R_Z^{(P)} R_Y^{(P)} R_X^{(P)}\bigr)^T$} \\
\bottomrule
\end{tabular}
\caption{GNC Aerospace vs.\ CTCT conventions---same product order,
transposed elementary matrices, transposed result.}
\end{table}

\begin{warningbox}[Common pitfall]
Both conventions are right-hand positive and use the same matrix product
order $R_Z\,R_Y\,R_X$. The \emph{only} difference is the sign of $\sin$ in
the elementary matrices. Mixing conventions silently transposes the result,
which can appear as a sign error in yaw or roll and is notoriously hard to
debug.
\end{warningbox}

\subsubsection*{Quick Diagnostic}

Compute $R_Z(90^\circ)\,\hat{x}$ where $\hat{x}=[1,0,0]^T$:
\begin{itemize}
  \item \textbf{GNC Aerospace} (passive): result $= [0,\,-1,\,0]^T$
        $\;\Rightarrow\;$ the nav $\hat{x}$ has body coordinate $-\hat{y}$
        (nav-to-body).
  \item \textbf{CTCT} (active): result $= [0,\,+1,\,0]^T$
        $\;\Rightarrow\;$ the body $\hat{x}$ points along nav $+\hat{y}$
        (body-to-nav).
\end{itemize}

\begin{keybox}[This document's convention]
All matrices in this document use the \textbf{GNC Aerospace} (passive)
convention. Therefore $R_Z(\psi)\,R_Y(\theta)\,R_X(\phi) = C_n^b$.
The body-to-nav matrix is obtained via transpose: $C_b^n = (C_n^b)^T$.
\end{keybox}

% ──────────────────────────────────────────────────────────────────────────────
\subsection{Active vs.\ Passive --- Worked Example}
% ──────────────────────────────────────────────────────────────────────────────

To make the distinction concrete, consider a single $90^\circ$ yaw rotation
applied to a vector $\vec{v}=[1,0,0]^T$ (pointing along the $X$-axis).

\subsubsection*{Active (CTCT) --- ``Move the vector''}

The vector physically rotates within a \textbf{fixed} frame.  The coordinate
axes do not move; the arrow does.

\begin{equation}
  R_Z^{(\text{active})}(90^\circ)
  = \begin{pmatrix}
      \cos 90^\circ & -\sin 90^\circ & 0 \\
      \sin 90^\circ &  \cos 90^\circ & 0 \\
      0             &  0             & 1
    \end{pmatrix}
  = \begin{pmatrix} 0 & -1 & 0 \\ 1 & 0 & 0 \\ 0 & 0 & 1 \end{pmatrix}
\end{equation}

\begin{equation}
  \vec{v}' = R_Z^{(\text{active})}(90^\circ)\,\vec{v}
  = \begin{pmatrix} 0\\1\\0 \end{pmatrix}
\end{equation}

The arrow that was pointing along $+X$ now points along $+Y$.
The axes did not move.

\begin{verbatim}
  Y                        Y
  |                        |
  |                        ^ v' = [0,1,0]
  |                        |
  +----> v = [1,0,0]  ->  +--------
       X  (fixed)               X  (fixed)

  BEFORE                   AFTER (active / CTCT)
\end{verbatim}

\emph{Think of it as:} you are standing still, watching an arrow on the ground
rotate in front of you.

\subsubsection*{Passive (GNC Aerospace) --- ``Move the frame''}

The vector stays \textbf{fixed in space}.  The coordinate frame rotates
underneath it, and we re-read the vector's coordinates in the new frame.

\begin{equation}
  R_Z^{(\text{passive})}(90^\circ)
  = \begin{pmatrix}
      \cos 90^\circ &  \sin 90^\circ & 0 \\
     -\sin 90^\circ &  \cos 90^\circ & 0 \\
      0             &  0             & 1
    \end{pmatrix}
  = \begin{pmatrix} 0 & 1 & 0 \\ -1 & 0 & 0 \\ 0 & 0 & 1 \end{pmatrix}
\end{equation}

\begin{equation}
  \vec{v}_{\text{new}} = R_Z^{(\text{passive})}(90^\circ)\,\vec{v}
  = \begin{pmatrix} 0\\-1\\0 \end{pmatrix}
\end{equation}

The arrow is still pointing the same direction in space, but we rotated the
coordinate frame $90^\circ$ CCW.  In this new frame the original arrow is at
$[0,-1,0]^T$---it points along $-Y'$.

\begin{verbatim}
  Y                        X'
  |                        |
  |                        |
  |                        |
  +----> v = [1,0,0]  ->  ---+--------
       X  (fixed)             |        Y'
                              |
                              v  v_new = [0,-1,0]

  BEFORE                   AFTER (passive / GNC Aerospace)
  (old frame)              (new frame rotated 90 deg CCW)
\end{verbatim}

\emph{Think of it as:} the arrow is nailed to the floor.  You rotate yourself
$90^\circ$ CCW.  From your new perspective the arrow that was ``forward'' is
now to your right---it reads $[0,-1,0]^T$ in your rotated frame.

\subsubsection*{Key Relationship}

The two matrices are \textbf{transposes} of each other for the same physical
rotation:

\begin{equation}
\boxed{
  R^{(\text{passive})} = \bigl(R^{(\text{active})}\bigr)^T
}
\end{equation}

Verify:
\[
  \begin{pmatrix} 0&1&0\\-1&0&0\\0&0&1 \end{pmatrix}
  = \begin{pmatrix} 0&-1&0\\1&0&0\\0&0&1 \end{pmatrix}^T
  \quad\checkmark
\]

\subsubsection*{3-D GNC Example}

An aircraft has yaw $\psi=90^\circ$ (nose pointing East).  A GPS gives
velocity in the nav frame: $\vec{v}_{\text{nav}} = [5,0,0]^T$\,m/s (moving
North).

\begin{itemize}
  \item \textbf{CTCT (Active)} --- ``Where does the body $\hat{x}$ point in
        nav?''
    \[
      R_Z^{(A)}(90^\circ)\begin{pmatrix}1\\0\\0\end{pmatrix}
      = \begin{pmatrix}0\\1\\0\end{pmatrix}
    \]
    Body $\hat{x}$ points along nav $+Y$ (East).\quad This is $C_b^n$.

  \item \textbf{GNC Aerospace (Passive)} --- ``What is the nav velocity in
        body coords?''
    \[
      R_Z^{(P)}(90^\circ)\begin{pmatrix}5\\0\\0\end{pmatrix}
      = \begin{pmatrix}0\\-5\\0\end{pmatrix}
    \]
    The 5\,m/s North velocity appears as $-5$\,m/s along body $Y$ (coming
    from the right).\quad This is $C_n^b$.
\end{itemize}

\subsubsection*{Summary}

\begin{table}[h]
\centering
\begin{tabular}{@{}p{2.5cm}p{5.5cm}p{5.5cm}@{}}
\toprule
 & \textbf{CTCT (Active/Alibi)} & \textbf{GNC Aerospace (Passive/Alias)} \\
\midrule
What moves?   & The vector             & The coordinate frame \\
What stays?   & The frame              & The vector \\
Result        & $\vec{v}'=R^{(A)}\vec{v}$ &
                $\vec{v}_{\text{new}}=R^{(P)}\vec{v}_{\text{old}}$ \\
GNC use       & $C_b^n$ (body-to-nav)  & $C_n^b$ (nav-to-body) \\
Relationship  & \multicolumn{2}{c}{$R^{(P)} = (R^{(A)})^T$ --- same physics,
                different bookkeeping} \\
\bottomrule
\end{tabular}
\caption{Active vs.\ passive rotation at a glance.}
\end{table}

% ──────────────────────────────────────────────────────────────────────────────
\subsection{Intrinsic/Extrinsic vs.\ Active/Passive --- Two Independent Choices}
% ──────────────────────────────────────────────────────────────────────────────

These two pairs of terms are often confused, but they answer
\textbf{different, independent questions}:

\begin{table}[h]
\centering
\begin{tabular}{@{}p{4.5cm}p{4.5cm}p{4.5cm}@{}}
\toprule
 & \textbf{What it decides} & \textbf{Options} \\
\midrule
Intrinsic vs.\ Extrinsic &
  Which \textbf{axes} to rotate about &
  Body-fixed (moving) vs.\ World-fixed (stationary) \\
Active vs.\ Passive &
  What the matrix \textbf{does} to a vector &
  Rotates the vector vs.\ Re-expresses it in a new frame \\
\bottomrule
\end{tabular}
\caption{The two independent convention choices.}
\end{table}

\subsubsection*{Intrinsic vs.\ Extrinsic --- ``About which axes?''}

This determines the \textbf{sequence of axes}:

\begin{itemize}
  \item \textbf{Intrinsic (body-fixed):} Each rotation is about the
        \emph{current} body axis, which moves with each step.\\
        \emph{``Yaw about body Z, then pitch about the new $Y'$, then roll
        about the new $X''$.''}
  \item \textbf{Extrinsic (world-fixed):} Each rotation is about the
        \emph{original} fixed world axis.\\
        \emph{``Rotate about fixed X, then fixed Y, then fixed Z.''}
\end{itemize}

Key fact: Intrinsic $Z$-$Y$-$X$ $=$ Extrinsic $X$-$Y$-$Z$ (reversed axis
order, same matrix product $R_Z\,R_Y\,R_X$).  This choice affects
\textbf{only the multiplication order}, nothing else.

\subsubsection*{Active vs.\ Passive --- ``What does the matrix do?''}

This determines the \textbf{meaning of the result}:

\begin{itemize}
  \item \textbf{Active (alibi):} The matrix physically \emph{moves} the
        vector to a new position within the same coordinate frame.
        \[
          \vec{v}' = R\,\vec{v}
          \quad\text{--- the vector rotates, the frame stays.}
        \]
  \item \textbf{Passive (alias):} The matrix \emph{re-expresses} the same
        (stationary) vector in a different coordinate frame.
        \[
          \vec{v}_{\text{new frame}} = R\,\vec{v}_{\text{old frame}}
          \quad\text{--- the frame rotates, the vector stays.}
        \]
\end{itemize}

This choice affects \textbf{only the sign of $\sin$} in the elementary
matrices (active and passive are transposes of each other).

% ──────────────────────────────────────────────────────────────────────────────
\subsection{Intrinsic/Extrinsic for GNC Aerospace vs.\ CTCT}
% ──────────────────────────────────────────────────────────────────────────────

Both GNC Aerospace and CTCT use the \textbf{same} axis sequence
(intrinsic $Z$-$Y$-$X$) and the \textbf{same} matrix product
$R_Z\,R_Y\,R_X$.  The only naming difference is which physical quantity
is assigned to the Y-axis and X-axis rotations.

\begin{table}[h]
\centering
\begin{tabular}{@{}p{3.5cm}p{5cm}p{5cm}@{}}
\toprule
 & \textbf{GNC Aerospace} & \textbf{CTCT} \\
\midrule
Axis sequence &
  Intrinsic $Z$-$Y$-$X$ &
  Intrinsic $Z$-$Y$-$X$ \textbf{(same)} \\
Extrinsic equivalent &
  Extrinsic $X$-$Y$-$Z$ &
  Extrinsic $X$-$Y$-$Z$ \textbf{(same)} \\
Product form &
  $R_Z\,R_Y\,R_X$ &
  $R_Z\,R_Y\,R_X$ \textbf{(same)} \\
\midrule
\multicolumn{3}{@{}l}{\textbf{Angle--axis assignment (where they differ):}} \\[4pt]
Z-axis rotation &
  Yaw ($\psi$) &
  Yaw ($\psi$) \textbf{(same)} \\
Y-axis rotation &
  Pitch ($\theta$) &
  Roll ($\phi_R$) \textbf{(swapped)} \\
X-axis rotation &
  Roll ($\phi$) &
  Pitch ($\theta_P$) \textbf{(swapped)} \\
\midrule
Intrinsic order &
  Yaw $\to$ Pitch $\to$ Roll &
  Yaw $\to$ Roll $\to$ Pitch \\
Extrinsic order &
  Roll $\to$ Pitch $\to$ Yaw &
  Pitch $\to$ Roll $\to$ Yaw \\
\bottomrule
\end{tabular}
\caption{Intrinsic/extrinsic comparison---same axis sequence
$Z$-$Y$-$X$ and same matrix product, but different angle names for
the Y and X axes.}
\end{table}

The intrinsic/extrinsic choice is \textbf{not} where the two conventions
differ.  They differ in (1)~active vs.\ passive and
(2)~which physical angle (pitch or roll) is assigned to the Y and X axes.

% ──────────────────────────────────────────────────────────────────────────────
\subsection{Active/Passive for GNC Aerospace vs.\ CTCT}
% ──────────────────────────────────────────────────────────────────────────────

Here is where the conventions diverge:

\begin{table}[h]
\centering
\begin{tabular}{@{}p{3.8cm}p{5cm}p{5cm}@{}}
\toprule
 & \textbf{GNC Aerospace} & \textbf{CTCT} \\
\midrule
Interpretation &
  Passive (alias) &
  Active (alibi) \\
$R_Z$ sign pattern &
  $+\sin$ above diagonal (row 1, col 2) &
  $+\sin$ below diagonal (row 2, col 1) \\
Product $R_Z R_Y R_X$ yields &
  $C_n^b$ (nav-to-body) &
  $C_b^n$ (body-to-nav) \\
What the matrix does &
  Re-expresses vector in body frame &
  Moves body vector into nav frame \\
To get the other direction &
  Transpose &
  Transpose \\
\bottomrule
\end{tabular}
\caption{Active/passive comparison---the \textbf{only} difference between
GNC Aerospace and CTCT.  Same angles, same order, transposed result.}
\end{table}

\begin{keybox}[Summary]
GNC Aerospace and CTCT use the \emph{same} rotation sequence (intrinsic
$Z$-$Y$-$X$) and the \emph{same} product order ($R_Z\,R_Y\,R_X$).
They differ \emph{only} in active vs.\ passive: each elementary matrix
is the transpose of the other convention's, and the composed result is
the transpose.

\[
  \underbrace{R_Z^{(A)} R_Y^{(A)} R_X^{(A)}}_{\text{CTCT: }C_b^n}
  \;=\;
  \underbrace{\bigl(R_Z^{(P)} R_Y^{(P)} R_X^{(P)}\bigr)^T}_{\text{Transpose of Aerospace: }(C_n^b)^T}
\]
\end{keybox}

\begin{warningbox}[CTCT Pitch/Roll Name Swap]
In CTCT code the Y-axis angle is called ``Roll'' and the X-axis angle
is called ``Pitch''---\textbf{swapped} relative to standard aerospace.
This applies to \texttt{Rotation\_Matrix.m},
\texttt{Calc\_Rotation\_NF\_Pltfrm1}, \texttt{EulerXYZBasis}, and most
autocode.  Newer C++ components (\texttt{DozerMinMaxYawCalibration})
use standard aerospace naming.  Always check which naming convention
a given file uses.
\end{warningbox}

\begin{keybox}[This document's choices]
\textbf{Intrinsic ZYX} (body-fixed axes) $+$ \textbf{Passive} (alias)
$\;\Longrightarrow\; R_Z\,R_Y\,R_X = C_n^b$ (nav-to-body).
\end{keybox}

% ══════════════════════════════════════════════════════════════════════════════
\section{Gimbal Lock}
% ══════════════════════════════════════════════════════════════════════════════

% ──────────────────────────────────────────────────────────────────────────────
\subsection{Definition}
% ──────────────────────────────────────────────────────────────────────────────

Gimbal lock occurs when two of the three rotation axes align, causing the
system to \textbf{lose one degree of freedom}. The representation collapses
from 3 independent rotations to effectively 2.

% ──────────────────────────────────────────────────────────────────────────────
\subsection{When Does It Happen?}
% ──────────────────────────────────────────────────────────────────────────────

For the intrinsic ZYX sequence, gimbal lock occurs when the \textbf{middle
angle} (pitch) reaches

\begin{equation}
  \theta = \pm\,90^\circ
\end{equation}

At $\theta = 90^\circ$ the Z and X axes become aligned, and yaw and roll
become indistinguishable---only their \emph{sum or difference} matters.

% ──────────────────────────────────────────────────────────────────────────────
\subsection{Axis Alignment at the Singularity}
% ──────────────────────────────────────────────────────────────────────────────

Starting with body frame $=$ nav frame (Z up, X forward, Y right):

\subsubsection*{Case 1: $\theta = +90^\circ$ (Nose Straight Up)}

After pitching by $+90^\circ$ about $Y'$:

\[
R_Y(90^\circ) = \begin{pmatrix} 0 & 0 & 1 \\ 0 & 1 & 0 \\ -1 & 0 & 0 \end{pmatrix}
\]

\begin{table}[h]
\centering
\begin{tabular}{@{}lll@{}}
\toprule
\textbf{Body Axis} & \textbf{Points Along (Nav)} & \textbf{Physical Meaning} \\
\midrule
$X''$ (forward) & $+Z_{\text{nav}}$ (up)       & Nose pointing straight up \\
$Y''$ (right)   & $+Y_{\text{nav}}$            & Unchanged \\
$Z''$ (belly)   & $-X_{\text{nav}}$ (backward) & Belly faces backward \\
\bottomrule
\end{tabular}
\end{table}

\begin{warningbox}[Gimbal Lock]
The body $X''$ axis is now aligned with the nav $Z$ axis:
\[
  X''_{\text{body}} \;\|\; Z_{\text{nav}}
\]
Roll (about $X''$) and yaw (about $Z_{\text{nav}}$) now rotate about the
\textbf{same axis}. One degree of freedom is lost.
\end{warningbox}

\subsubsection*{Case 2: $\theta = -90^\circ$ (Nose Straight Down)}

\begin{table}[h]
\centering
\begin{tabular}{@{}lll@{}}
\toprule
\textbf{Body Axis} & \textbf{Points Along (Nav)} & \textbf{Physical Meaning} \\
\midrule
$X''$ (forward) & $-Z_{\text{nav}}$ (down) & Nose pointing straight down \\
$Y''$ (right)   & $+Y_{\text{nav}}$        & Unchanged \\
$Z''$ (belly)   & $+X_{\text{nav}}$ (fwd)  & Belly faces forward \\
\bottomrule
\end{tabular}
\end{table}

Again $X'' \| Z_{\text{nav}}$ (anti-parallel). Now $(\psi + \phi)$ is the
only observable combined parameter.

% ──────────────────────────────────────────────────────────────────────────────
\subsection{Proof: Why Exactly One DOF Is Lost}
% ──────────────────────────────────────────────────────────────────────────────

Expand the full rotation matrix at $\theta = 90^\circ$:

\begin{equation}
  R = R_Z(\psi)\; R_Y(90^\circ)\; R_X(\phi)
    = \begin{pmatrix}
        0 & \sin(\psi - \phi) & \cos(\psi - \phi) \\
        0 & \cos(\psi - \phi) & -\sin(\psi - \phi) \\
       -1 & 0                 & 0
      \end{pmatrix}
  \label{eq:gimbal}
\end{equation}

\textbf{Every entry depends only on} $(\psi - \phi)$, never on $\psi$ or
$\phi$ individually. Therefore infinitely many $(\psi,\phi)$ pairs yield the
same orientation:

\begin{table}[h]
\centering
\begin{tabular}{@{}cccc@{}}
\toprule
$\psi$ (yaw) & $\phi$ (roll) & $\psi - \phi$ & Same matrix? \\
\midrule
$30^\circ$  & $0^\circ$    & $30^\circ$ & \checkmark \\
$40^\circ$  & $10^\circ$   & $30^\circ$ & \checkmark \\
$0^\circ$   & $-30^\circ$  & $30^\circ$ & \checkmark \\
$100^\circ$ & $70^\circ$   & $30^\circ$ & \checkmark \\
\bottomrule
\end{tabular}
\caption{Multiple Euler angle triples mapping to the same orientation at
$\theta = 90^\circ$.}
\end{table}

You have 3 parameters $(\psi, \theta, \phi)$ but only \textbf{2} of them
independently change the orientation:
\begin{enumerate}
  \item $\theta$ --- but it is stuck at $90^\circ$.
  \item $(\psi - \phi)$ --- one combined number, not two.
\end{enumerate}

That is \textbf{2 effective DOF}, not 3. One is lost.

% ──────────────────────────────────────────────────────────────────────────────
\subsection{Analogy: Latitude and Longitude at the Pole}
% ──────────────────────────────────────────────────────────────────────────────

Consider a map with latitude and longitude:

\begin{itemize}
  \item At the \textbf{equator}, changing latitude moves you north/south and
        changing longitude moves you east/west.
        \textbf{Two independent directions.}
  \item At the \textbf{North Pole}, changing longitude does nothing---every
        longitude passes through the same point. Longitude is meaningless.
        \textbf{One independent direction only.}
\end{itemize}

The Euler angle singularity is the same idea. At $\theta = 90^\circ$, $\psi$
and $\phi$ collapse into a single effective parameter $(\psi - \phi)$, just as
at the pole all longitudes collapse to one point.

% ──────────────────────────────────────────────────────────────────────────────
\subsection{Consequences}
% ──────────────────────────────────────────────────────────────────────────────

\begin{table}[h]
\centering
\begin{tabular}{@{}p{4cm}p{9cm}@{}}
\toprule
\textbf{Effect} & \textbf{Description} \\
\midrule
Ambiguous decomposition & Infinitely many $(\psi,\phi)$ pairs produce the
  same rotation. \\
Numerical instability & \texttt{atan2}-based extraction becomes
  ill-conditioned near $\theta = \pm 90^\circ$. \\
Discontinuous outputs & Small physical rotations produce large jumps in
  extracted angles. \\
Not a physical limitation & The orientation itself is fine; only the
  \emph{Euler angle representation} breaks. \\
\bottomrule
\end{tabular}
\caption{Effects of gimbal lock.}
\end{table}

% ──────────────────────────────────────────────────────────────────────────────
\subsection{What ``Losing a DOF'' Does and Does Not Mean}
% ──────────────────────────────────────────────────────────────────────────────

\textbf{It does NOT mean:}
\begin{itemize}[label=$\times$]
  \item The body cannot physically reach certain orientations.
  \item The body is stuck or constrained.
  \item There are only 2 possible directions to move.
\end{itemize}

\textbf{It DOES mean:}
\begin{itemize}[label=\checkmark]
  \item The \emph{representation} (three Euler angles) can no longer uniquely
        describe orientations near $\theta = 90^\circ$.
  \item Two of the three angle ``knobs'' ($\psi$ and $\phi$) turn the same
        physical axis.
  \item Small physical rotations about one axis produce wild swings or
        ambiguity in the extracted angles.
\end{itemize}

% ──────────────────────────────────────────────────────────────────────────────
\subsection{Detection}
% ──────────────────────────────────────────────────────────────────────────────

\begin{lstlisting}[style=matlab, caption={Detecting gimbal lock in MATLAB.}]
% After extracting pitch (theta) from a rotation matrix:
if abs(cos(theta)) < 1e-6
    warning('Near gimbal lock -- Euler angles unreliable');
end
\end{lstlisting}

% ──────────────────────────────────────────────────────────────────────────────
\subsection{Mitigation Strategies}
% ──────────────────────────────────────────────────────────────────────────────

\begin{table}[h]
\centering
\begin{tabular}{@{}p{4cm}p{9cm}@{}}
\toprule
\textbf{Approach} & \textbf{Trade-off} \\
\midrule
Quaternions & No singularities, 4 parameters, must normalise.
  Standard in GNC/INS. \\
Rotation matrices (DCM) & No singularities, 9 parameters (6 redundant),
  must re-orthogonalise. \\
Different Euler order & Moves the singularity (e.g.\ ZYZ puts it at
  $\theta = 0^\circ$), does not eliminate it. \\
Clamp \& flag & Keep Euler angles for display but clamp $\theta$
  and flag the degenerate case. \\
\bottomrule
\end{tabular}
\caption{Strategies to avoid or mitigate gimbal lock.}
\end{table}

\begin{keybox}[Key Takeaway]
Gimbal lock is an inherent limitation of \textbf{all} 3-parameter attitude
representations (Euler angles, Rodrigues parameters). It cannot be
``fixed''---only avoided by using a higher-dimensional representation
(quaternions or DCMs) for computation, and reserving Euler angles for
human-readable output where $\theta$ stays well away from $\pm 90^\circ$.
\end{keybox}

% ══════════════════════════════════════════════════════════════════════════════
\section{Intrinsic vs Extrinsic: Axis Alignment Comparison}
% ══════════════════════════════════════════════════════════════════════════════

\begin{table}[h]
\centering
\begin{tabular}{@{}p{3.5cm}p{5.5cm}p{5.5cm}@{}}
\toprule
 & \textbf{Intrinsic ZYX} & \textbf{Extrinsic XYZ} \\
\midrule
What aligns? & Body $X''$ aligns with nav $Z$ &
  After fixed-$Y$ rotation of $90^\circ$, fixed-$Z$ rotation acts the same as
  the earlier fixed-$X$ rotation \\
Locked axes & Roll (body $X''$) $\leftrightarrow$ Yaw (nav $Z$) &
  First rotation (fixed $X$) $\leftrightarrow$ Third rotation (fixed $Z$) \\
Lost DOF & Only $(\psi \pm \phi)$ observable &
  Only $(\alpha \pm \gamma)$ observable \\
Intuition & Easier: ``two body axes collapse onto same nav axis'' &
  Harder: ``fixed-frame rotations become redundant'' \\
\bottomrule
\end{tabular}
\caption{Gimbal lock comparison between intrinsic and extrinsic conventions.}
\end{table}

% ══════════════════════════════════════════════════════════════════════════════
\section{MATLAB Example}
% ══════════════════════════════════════════════════════════════════════════════

\begin{lstlisting}[style=matlab, caption={Intrinsic ZYX and equivalent
extrinsic XYZ in MATLAB.}]
% Intrinsic ZYX (yaw=30 deg, pitch=20 deg, roll=10 deg)
R_intrinsic = rotz(30) * roty(20) * rotx(10);

% Equivalent extrinsic XYZ (same angles, reversed order)
% The matrix product is identical:
R_extrinsic = rotz(30) * roty(20) * rotx(10);

% Verify
assert(norm(R_intrinsic - R_extrinsic) < 1e-12);

% Nav-to-body
v_body = R_intrinsic * v_nav;

% Body-to-nav (transpose)
v_nav  = R_intrinsic' * v_body;
\end{lstlisting}

% ══════════════════════════════════════════════════════════════════════════════
\section{Quaternions}
% ══════════════════════════════════════════════════════════════════════════════

Quaternions provide a singularity-free, compact representation of 3D
rotations. They use four parameters with one constraint ($\|q\|=1$),
avoiding the gimbal lock problem inherent in all three-parameter
representations.

% ──────────────────────────────────────────────────────────────────────────────
\subsection{Definition}
% ──────────────────────────────────────────────────────────────────────────────

A quaternion is a four-dimensional hypercomplex number:

\begin{equation}
  q = w + x\,\mathbf{i} + y\,\mathbf{j} + z\,\mathbf{k}
\end{equation}

where $w,x,y,z \in \mathbb{R}$ and the imaginary units satisfy:

\begin{equation}
  \mathbf{i}^2 = \mathbf{j}^2 = \mathbf{k}^2
  = \mathbf{i}\,\mathbf{j}\,\mathbf{k} = -1
\end{equation}

A quaternion is often written as a scalar--vector pair:

\begin{equation}
  q = (w,\;\vec{v}), \quad \vec{v} = (x,\,y,\,z)
\end{equation}

% ──────────────────────────────────────────────────────────────────────────────
\subsection{Multiplication Rules}
% ──────────────────────────────────────────────────────────────────────────────

From the imaginary-unit identities the basis multiplication table follows:

\begin{table}[h]
\centering
\begin{tabular}{@{}c|ccc@{}}
\toprule
$\times$ & $\mathbf{i}$ & $\mathbf{j}$ & $\mathbf{k}$ \\
\midrule
$\mathbf{i}$ & $-1$ & $\mathbf{k}$  & $-\mathbf{j}$ \\
$\mathbf{j}$ & $-\mathbf{k}$ & $-1$ & $\mathbf{i}$  \\
$\mathbf{k}$ & $\mathbf{j}$  & $-\mathbf{i}$ & $-1$ \\
\bottomrule
\end{tabular}
\caption{Hamilton quaternion basis multiplication table.}
\end{table}

% ──────────────────────────────────────────────────────────────────────────────
\subsection{Quaternion Arithmetic}
% ──────────────────────────────────────────────────────────────────────────────

\textbf{Addition} (element-wise):
\begin{equation}
  q_1 + q_2 = (w_1+w_2,\;\vec{v}_1+\vec{v}_2)
\end{equation}

\textbf{Multiplication} (non-commutative):
\begin{equation}
  q_1\,q_2 = \bigl(w_1 w_2 - \vec{v}_1\cdot\vec{v}_2,\;
              w_1\vec{v}_2 + w_2\vec{v}_1
              + \vec{v}_1\times\vec{v}_2\bigr)
\end{equation}

The cross-product term $\vec{v}_1\times\vec{v}_2$ makes quaternion
multiplication \textbf{non-commutative}: $q_1 q_2 \neq q_2 q_1$ in general.

\textbf{Conjugate}:
\begin{equation}
  q^{*} = (w,\;-\vec{v})
\end{equation}

\textbf{Norm}:
\begin{equation}
  \|q\| = \sqrt{w^2 + x^2 + y^2 + z^2}
\end{equation}

\textbf{Inverse}:
\begin{equation}
  q^{-1} = \frac{q^{*}}{\|q\|^2}
\end{equation}

For \textbf{unit quaternions} ($\|q\|=1$), the inverse equals the conjugate:
$q^{-1} = q^{*}$.

% ──────────────────────────────────────────────────────────────────────────────
\subsection{Unit Quaternions as Rotations}
% ──────────────────────────────────────────────────────────────────────────────

A unit quaternion encodes a rotation about axis $\hat{n}$ by angle $\theta$:

\begin{equation}
  q = \left(\cos\frac{\theta}{2},\;
      \sin\frac{\theta}{2}\;\hat{n}\right)
\end{equation}

The \textbf{half-angle} arises from the double-cover property (Section
\ref{sec:double-cover}).

\subsubsection*{Rotating a Vector}

To rotate a vector $\vec{p}$, embed it as a pure quaternion
$p=(0,\vec{p})$, then compute:

\begin{equation}
  p' = q\,p\,q^{*}
\end{equation}

The vector part of $p'$ is the rotated vector.

% ──────────────────────────────────────────────────────────────────────────────
\subsection{Composing Rotations}
% ──────────────────────────────────────────────────────────────────────────────

Applying rotation $q_1$ first, then $q_2$, is a single quaternion
multiplication:

\begin{equation}
  q_{\text{total}} = q_2\,q_1
\end{equation}

The product of two unit quaternions is again a unit quaternion, so no
renormalisation is needed after composition (unlike numerical integration).

% ──────────────────────────────────────────────────────────────────────────────
\subsection{Conversions}
% ──────────────────────────────────────────────────────────────────────────────

\subsubsection*{Quaternion $\to$ Rotation Matrix}

\begin{equation}
  R = \begin{pmatrix}
    1-2(y^2+z^2) & 2(xy-wz)     & 2(xz+wy) \\
    2(xy+wz)     & 1-2(x^2+z^2) & 2(yz-wx) \\
    2(xz-wy)     & 2(yz+wx)     & 1-2(x^2+y^2)
  \end{pmatrix}
\end{equation}

\subsubsection*{Quaternion $\to$ Euler Angles (ZYX)}

\begin{align}
  \phi   &= \operatorname{atan2}\!\bigl(2(wx+yz),\;1-2(x^2+y^2)\bigr) \\
  \theta &= \arcsin\!\bigl(2(wy-zx)\bigr) \\
  \psi   &= \operatorname{atan2}\!\bigl(2(wz+xy),\;1-2(y^2+z^2)\bigr)
\end{align}

\subsubsection*{Axis-Angle $\to$ Quaternion}

\begin{equation}
  q = \left(\cos\frac{\theta}{2},\;
      \sin\frac{\theta}{2}\;\hat{n}\right)
\end{equation}

% ──────────────────────────────────────────────────────────────────────────────
\subsection{Double Cover of $SO(3)$}
\label{sec:double-cover}
% ──────────────────────────────────────────────────────────────────────────────

Unit quaternions live on $S^3$ (the 3-sphere in $\mathbb{R}^4$). The map
$q \mapsto R(q)$ from unit quaternions to the rotation group $SO(3)$ is a
\textbf{2-to-1 surjection}: both $q$ and $-q$ produce the same rotation
matrix.

\begin{equation}
  S^3 \big/ \{\pm 1\} \;\cong\; SO(3)
\end{equation}

This is why the half-angle appears: a full $2\pi$ physical rotation
corresponds to only a $\pi$ traversal on $S^3$. A $4\pi$ rotation is needed
to return to the \emph{same} quaternion.

\begin{keybox}[Practical consequence]
Because $q$ and $-q$ are the same rotation, algorithms that interpolate or
average quaternions must check the sign of $q_1 \cdot q_2$ and negate one
quaternion if needed to ensure the shortest-path interpretation.
\end{keybox}

% ──────────────────────────────────────────────────────────────────────────────
\subsection{Quaternion Derivative --- Rotational Kinematics}
% ──────────────────────────────────────────────────────────────────────────────

If a rigid body has angular velocity $\vec{\omega}$ expressed in the body
frame, the quaternion evolves as:

\begin{equation}
  \dot{q} = \tfrac{1}{2}\,q \otimes \omega_q,
  \qquad \omega_q = (0,\,\vec{\omega})
\end{equation}

In matrix form:

\begin{equation}
  \dot{q} = \frac{1}{2}\,\Omega\,q, \qquad
  \Omega = \begin{pmatrix}
    0         & -\omega_x & -\omega_y & -\omega_z \\
    \omega_x  & 0          & \omega_z  & -\omega_y \\
    \omega_y  & -\omega_z  & 0          & \omega_x  \\
    \omega_z  & \omega_y   & -\omega_x  & 0
  \end{pmatrix}
\end{equation}

This ODE is integrated in INS/AHRS systems to propagate attitude.

% ──────────────────────────────────────────────────────────────────────────────
\subsection{SLERP --- Spherical Linear Interpolation}
% ──────────────────────────────────────────────────────────────────────────────

To smoothly interpolate between orientations $q_0$ and $q_1$:

\begin{equation}
  \operatorname{slerp}(q_0,q_1,t)
  = q_0\,\frac{\sin\bigl((1-t)\,\Omega\bigr)}{\sin\Omega}
  + q_1\,\frac{\sin(t\,\Omega)}{\sin\Omega},
  \qquad t\in[0,1]
\end{equation}

where $\cos\Omega = q_0 \cdot q_1$ (4-D dot product). SLERP traverses the
great circle on $S^3$, producing constant angular velocity---unlike linear
interpolation of Euler angles.

% ──────────────────────────────────────────────────────────────────────────────
\subsection{Quaternion Error and Attitude Control}
% ──────────────────────────────────────────────────────────────────────────────

The rotation \emph{from} $q_a$ \emph{to} $q_b$ is:

\begin{equation}
  q_{\text{err}} = q_b \otimes q_a^{*}
\end{equation}

For a proportional attitude controller the error signal is commonly
extracted as:

\begin{equation}
  \vec{e} = 2\,\operatorname{sign}(w_e)\,\vec{v}_e,
  \qquad (w_e,\vec{v}_e) = q_{\text{err}}
\end{equation}

For small angles this approximates the rotation-vector error.

% ──────────────────────────────────────────────────────────────────────────────
\subsection{Exponential and Logarithmic Maps}
% ──────────────────────────────────────────────────────────────────────────────

These connect unit quaternions to the Lie algebra
$\mathfrak{so}(3)$.

\textbf{Exponential map} (rotation vector $\to$ quaternion):
\begin{equation}
  \exp(\vec{\theta})
  = \left(\cos\frac{\|\vec{\theta}\|}{2},\;
    \frac{\vec{\theta}}{\|\vec{\theta}\|}
    \sin\frac{\|\vec{\theta}\|}{2}\right)
\end{equation}

\textbf{Logarithmic map} (quaternion $\to$ rotation vector):
\begin{equation}
  \log(q)
  = 2\,\operatorname{atan2}\!\bigl(\|\vec{v}\|,w\bigr)\;
    \frac{\vec{v}}{\|\vec{v}\|}
\end{equation}

These maps are essential for:
\begin{itemize}
  \item \textbf{Optimisation on $SO(3)$}: perturb in the tangent space, then
        map back.
  \item \textbf{Kalman filtering}: attitude errors live in three-dimensional
        tangent space (multiplicative EKF / MEKF).
  \item \textbf{Averaging rotations}: compute the mean in log space.
\end{itemize}

% ──────────────────────────────────────────────────────────────────────────────
\subsection{Euler Angles vs.\ Quaternions}
% ──────────────────────────────────────────────────────────────────────────────

\begin{table}[h]
\centering
\begin{tabular}{@{}p{4cm}p{4.5cm}p{4.5cm}@{}}
\toprule
\textbf{Property} & \textbf{Euler Angles} & \textbf{Quaternions} \\
\midrule
Parameters        & 3                      & 4 (with 1 constraint) \\
Singularities     & Gimbal lock at $\pm90^\circ$ pitch & None \\
Composition       & Trigonometry-heavy     & Single multiplication \\
Interpolation     & Non-uniform, jerky     & SLERP (smooth) \\
Numerical drift   & N/A                    & Re-normalise: $q\leftarrow q/\|q\|$ \\
Human intuition   & Easy                   & Harder initially \\
\bottomrule
\end{tabular}
\caption{Comparison of Euler angles and quaternions.}
\end{table}

% ──────────────────────────────────────────────────────────────────────────────
\subsection{Practical Tips for GNC}
% ──────────────────────────────────────────────────────────────────────────────

\begin{enumerate}
  \item \textbf{Normalisation.} After every integration step, re-normalise:
        $q \leftarrow q/\|q\|$.
  \item \textbf{Convention.} Hamilton ($w$ first) vs.\ JPL ($w$ last) changes
        composition order. Know which your codebase uses.
  \item \textbf{Sign convention.} Enforce $w>0$ to avoid discontinuities in
        logging and control.
  \item \textbf{Small-angle approximation.}
        For $\theta \ll 1$: $q \approx (1,\,\vec{\theta}/2)$.
  \item \textbf{MEKF.} Represent attitude error as a 3-vector (rotation
        vector) to avoid the rank deficiency of estimating 4 quaternion
        components with a $4\times4$ covariance.
\end{enumerate}

% ──────────────────────────────────────────────────────────────────────────────
\subsection{MATLAB Examples}
% ──────────────────────────────────────────────────────────────────────────────

\begin{lstlisting}[style=matlab, caption={Quaternion rotation of a vector.}]
% Define a unit quaternion: 45-deg rotation about Z
theta = deg2rad(45);
q = [cos(theta/2), 0, 0, sin(theta/2)];   % [w x y z]

% Rotate a vector using quaternion sandwich product
v = [1; 0; 0];
v_q = [0, v'];                              % pure quaternion
v_rot = quatmultiply(quatmultiply(q, v_q), quatconj(q));
v_rotated = v_rot(2:4)'                     % extract vector part
\end{lstlisting}

\begin{lstlisting}[style=matlab, caption={Quaternion to Euler (ZYX) and back.}]
% quaternion -> Euler ZYX
eul = quat2eul(q, 'ZYX');   % [yaw pitch roll]

% Euler ZYX -> quaternion
q_back = eul2quat(eul, 'ZYX');
assert(norm(q - q_back) < 1e-12);
\end{lstlisting}

\begin{lstlisting}[style=matlab, caption={SLERP between two orientations.}]
q0 = [1 0 0 0];                     % identity
q1 = [cos(pi/4) 0 sin(pi/4) 0];     % 90 deg about Y
t  = 0.5;
q_mid = quatinterp(q0, q1, t, 'slerp');
\end{lstlisting}

% ======================================================================
\section{Denavit--Hartenberg (DH) Parameters}
% ======================================================================

The Denavit--Hartenberg convention provides a systematic way to describe
the geometry of a serial kinematic chain (robot arm, excavator boom,
crane, etc.)\ using only \textbf{four parameters per joint}.  Combined
with a simple $4\times4$ homogeneous transformation, DH parameters let
you compute the position and orientation of the end effector from the
joint angles alone.

% ----------------------------------------------------------------------
\subsection{The Four DH Parameters}
% ----------------------------------------------------------------------

Each link $i$ in the chain is described by exactly four quantities:

\begin{table}[h]
\centering
\begin{tabular}{@{}clp{8.5cm}@{}}
\toprule
\textbf{Symbol} & \textbf{Name} & \textbf{Definition} \\
\midrule
$\theta_i$ & Joint angle &
  Angle about $Z_{i-1}$ from $X_{i-1}$ to $X_i$
  (measured in the $X_{i-1}$--$Y_{i-1}$ plane). \\[4pt]
$d_i$      & Link offset &
  Distance along $Z_{i-1}$ from the origin of frame $i{-}1$
  to the intersection of $Z_{i-1}$ and $X_i$. \\[4pt]
$a_i$      & Link length &
  Distance along $X_i$ from $Z_{i-1}$ to $Z_i$
  (the ``common normal'' length). \\[4pt]
$\alpha_i$ & Link twist &
  Angle about $X_i$ from $Z_{i-1}$ to $Z_i$
  (measured in the $Y_i$--$Z_i$ plane). \\
\bottomrule
\end{tabular}
\caption{The four Denavit--Hartenberg parameters for link $i$.}
\end{table}

\begin{keybox}[DH frame assignment rules]
\begin{enumerate}
  \item The $Z_{i-1}$ axis lies along the axis of joint $i$.
  \item The $X_i$ axis is the common normal from $Z_{i-1}$ to $Z_i$,
        pointing away from $Z_{i-1}$.
  \item The $Y_i$ axis completes the right-hand frame:
        $Y_i = Z_i \times X_i$.
  \item The origin of frame $i$ is at the intersection of $X_i$ and
        $Z_i$ (or at the foot of the common normal on $Z_i$).
\end{enumerate}
\end{keybox}

% ----------------------------------------------------------------------
\subsection{The DH Transformation Matrix}
% ----------------------------------------------------------------------

The homogeneous transformation from frame $i{-}1$ to frame $i$ is:

\begin{equation}
\boxed{
  {}^{i-1}T_i =
  \underbrace{
    \operatorname{Rot}_{Z}(\theta_i)\;
    \operatorname{Trans}_{Z}(d_i)
  }_{\text{along }Z_{i-1}}
  \;
  \underbrace{
    \operatorname{Trans}_{X}(a_i)\;
    \operatorname{Rot}_{X}(\alpha_i)
  }_{\text{along }X_i}
}
\label{eq:DH-decomp}
\end{equation}

Multiplying these four elementary transforms out:

\begin{equation}
\boxed{
  {}^{i-1}T_i =
  \begin{pmatrix}
    \cos\theta_i & -\sin\theta_i\cos\alpha_i &
      \sin\theta_i\sin\alpha_i  & a_i\cos\theta_i \\[3pt]
    \sin\theta_i &  \cos\theta_i\cos\alpha_i &
      -\cos\theta_i\sin\alpha_i & a_i\sin\theta_i \\[3pt]
    0 & \sin\alpha_i & \cos\alpha_i & d_i \\[3pt]
    0 & 0 & 0 & 1
  \end{pmatrix}
}
\label{eq:DH-matrix}
\end{equation}

The upper-left $3\times3$ block is the rotation ${}^{i-1}R_i$; the
right-most column (top three entries) is the translation
${}^{i-1}\vec{p}_i$.

\begin{keybox}[Reading the matrix]
\begin{itemize}
  \item \textbf{Row 3, col 4}: $d_i$ --- the offset along the previous
        joint axis.
  \item \textbf{Rows 1--2, col 4}: $a_i\cos\theta_i$ and
        $a_i\sin\theta_i$ --- the link-length contribution projected
        by the joint angle.
  \item \textbf{Column 3}: contains $\sin\alpha_i$, $\cos\alpha_i$ ---
        the twist between successive $Z$ axes.
\end{itemize}
\end{keybox}

% ----------------------------------------------------------------------
\subsection{Rotary (Revolute) vs.\ Linear (Prismatic) Joints}
% ----------------------------------------------------------------------

The four DH parameters have different roles depending on the joint type:

\begin{table}[h]
\centering
\begin{tabular}{@{}p{3.2cm}p{5.5cm}p{5.5cm}@{}}
\toprule
 & \textbf{Revolute (Rotary) Joint} &
   \textbf{Prismatic (Linear) Joint} \\
\midrule
Joint variable &
  $\theta_i$ (angle) &
  $d_i$ (displacement) \\[4pt]
Fixed parameters &
  $d_i,\; a_i,\; \alpha_i$ &
  $\theta_i,\; a_i,\; \alpha_i$ \\[4pt]
Motion &
  Rotation about $Z_{i-1}$ &
  Translation along $Z_{i-1}$ \\[4pt]
DH matrix changes &
  $\cos\theta_i$, $\sin\theta_i$ vary &
  $d_i$ varies (row 3, col 4) \\[4pt]
Typical use &
  Robot elbow, excavator boom pivot,
  steering joint &
  Hydraulic cylinder, telescoping mast,
  linear actuator \\
\bottomrule
\end{tabular}
\caption{DH parameter roles for revolute vs.\ prismatic joints.}
\end{table}

\begin{warningbox}[Identifying the joint variable]
For \textbf{any} joint, exactly \textbf{one} of the four DH parameters
is the joint variable; the other three are fixed by the mechanism
geometry.
\begin{itemize}
  \item Revolute: the joint variable is $\theta_i$.
  \item Prismatic: the joint variable is $d_i$.
\end{itemize}
All higher-order joints (helical, cylindrical, spherical) can be
decomposed into combinations of these two.
\end{warningbox}

\subsubsection*{Revolute Joint --- Details}

For a revolute joint at link $i$:
\begin{itemize}
  \item $\theta_i$ is the \textbf{unknown} (measured by encoder or
        commanded by controller).
  \item $d_i$ is a fixed offset (often zero for planar mechanisms).
  \item $a_i$ is the physical link length (boom segment, forearm, etc.).
  \item $\alpha_i$ is the fixed twist angle between consecutive axes
        (e.g.\ $0^\circ$ for parallel axes, $90^\circ$ for
        perpendicular).
\end{itemize}

\subsubsection*{Prismatic Joint --- Details}

For a prismatic joint at link $i$:
\begin{itemize}
  \item $d_i$ is the \textbf{unknown} (cylinder extension).
  \item $\theta_i$ is a fixed constant (often $0^\circ$).
  \item $a_i$ is typically $0$ (the slide direction \emph{is} the
        $Z_{i-1}$ axis).
  \item $\alpha_i$ sets the twist to the next joint's axis.
\end{itemize}

% ----------------------------------------------------------------------
\subsection{Forward Kinematics --- End Effector Position}
% ----------------------------------------------------------------------

For an $n$-joint serial chain the end effector pose (position $+$
orientation) in the base frame is:

\begin{equation}
\boxed{
  {}^{0}T_n = {}^{0}T_1 \;\; {}^{1}T_2 \;\; {}^{2}T_3 \;\cdots\;
  {}^{n-1}T_n
  = \prod_{i=1}^{n} {}^{i-1}T_i
}
\label{eq:FK}
\end{equation}

Each ${}^{i-1}T_i$ is the $4\times4$ DH matrix from
Equation~\eqref{eq:DH-matrix}, evaluated with the current joint
variable ($\theta_i$ or $d_i$).

From the resulting $4\times4$ matrix:

\begin{equation}
  {}^{0}T_n =
  \begin{pmatrix}
    {}^{0}R_n & {}^{0}\vec{p}_n \\
    \vec{0}^T & 1
  \end{pmatrix}
\end{equation}

\begin{itemize}
  \item ${}^{0}\vec{p}_n = \begin{pmatrix} p_x \\ p_y \\ p_z \end{pmatrix}$
        is the \textbf{end effector position} in the base frame.
  \item ${}^{0}R_n$ is the \textbf{end effector orientation} (a $3\times3$
        rotation matrix that can be converted to Euler angles or a
        quaternion using the formulas in earlier sections).
\end{itemize}

\subsubsection*{Worked Example: 2-Link Planar Arm (Two Revolute Joints)}

\begin{table}[h]
\centering
\begin{tabular}{@{}ccccc@{}}
\toprule
\textbf{Link} & $\theta_i$ & $d_i$ & $a_i$ & $\alpha_i$ \\
\midrule
1 & $\theta_1$ (variable) & 0 & $L_1$ & $0^\circ$ \\
2 & $\theta_2$ (variable) & 0 & $L_2$ & $0^\circ$ \\
\bottomrule
\end{tabular}
\caption{DH table for a 2-DOF planar arm.}
\end{table}

With $\alpha_i = 0$ and $d_i = 0$, each DH matrix simplifies to:

\begin{equation}
  {}^{i-1}T_i =
  \begin{pmatrix}
    \cos\theta_i & -\sin\theta_i & 0 & L_i\cos\theta_i \\
    \sin\theta_i &  \cos\theta_i & 0 & L_i\sin\theta_i \\
    0 & 0 & 1 & 0 \\
    0 & 0 & 0 & 1
  \end{pmatrix}
\end{equation}

The end effector position is:

\begin{equation}
\boxed{
  \begin{pmatrix} p_x \\ p_y \end{pmatrix}
  = \begin{pmatrix}
      L_1\cos\theta_1 + L_2\cos(\theta_1+\theta_2) \\
      L_1\sin\theta_1 + L_2\sin(\theta_1+\theta_2)
    \end{pmatrix}
}
\end{equation}

This is the classic 2-link forward kinematics result.

\subsubsection*{Worked Example: Revolute + Prismatic (RP)}

Consider a rotary base followed by a telescoping arm:

\begin{table}[h]
\centering
\begin{tabular}{@{}ccccc@{}}
\toprule
\textbf{Link} & $\theta_i$ & $d_i$ & $a_i$ & $\alpha_i$ \\
\midrule
1 (revolute)  & $\theta_1$ (variable) & 0   & 0 & $90^\circ$ \\
2 (prismatic) & $0^\circ$             & $d_2$ (variable) & 0 & $0^\circ$ \\
\bottomrule
\end{tabular}
\caption{DH table for a revolute--prismatic (RP) mechanism.}
\end{table}

Joint 1 rotates about the base $Z_0$; joint 2 extends along $Z_1$
(perpendicular to $Z_0$ due to $\alpha_1 = 90^\circ$).

\begin{equation}
  {}^{0}T_1 =
  \begin{pmatrix}
    \cos\theta_1 & 0 & \sin\theta_1 & 0 \\
    \sin\theta_1 & 0 & -\cos\theta_1 & 0 \\
    0 & 1 & 0 & 0 \\
    0 & 0 & 0 & 1
  \end{pmatrix},\quad
  {}^{1}T_2 =
  \begin{pmatrix}
    1 & 0 & 0 & 0 \\
    0 & 1 & 0 & 0 \\
    0 & 0 & 1 & d_2 \\
    0 & 0 & 0 & 1
  \end{pmatrix}
\end{equation}

End effector position:

\begin{equation}
\boxed{
  {}^{0}\vec{p}_2 =
  \begin{pmatrix}
    d_2\,\sin\theta_1 \\ -d_2\,\cos\theta_1 \\ 0
  \end{pmatrix}
}
\end{equation}

The prismatic extension $d_2$ is projected through the base rotation
$\theta_1$---exactly how a turret-mounted telescoping arm behaves.

% ----------------------------------------------------------------------
\subsection{Summary of DH Procedure}
% ----------------------------------------------------------------------

\begin{enumerate}
  \item \textbf{Number the joints} 1 to $n$ from base to end effector.
  \item \textbf{Assign $Z$ axes}: $Z_{i-1}$ along the axis of joint $i$.
  \item \textbf{Assign $X$ axes}: $X_i$ along the common normal from
        $Z_{i-1}$ to $Z_i$.
  \item \textbf{Fill the DH table}: for each link, record
        $\theta_i$, $d_i$, $a_i$, $\alpha_i$.
  \item \textbf{Mark the joint variable}: $\theta_i$ for revolute,
        $d_i$ for prismatic.
  \item \textbf{Build ${}^{i-1}T_i$}: plug each row of the table into
        Equation~\eqref{eq:DH-matrix}.
  \item \textbf{Multiply}: ${}^{0}T_n = \prod {}^{i-1}T_i$ to get the
        end effector pose.
\end{enumerate}

\begin{keybox}[DH vs.\ Euler angles]
DH parameters describe the \emph{geometry between successive joints}
(link lengths, offsets, twists), while Euler angles describe the
\emph{orientation of a single body}.  A serial chain's end effector
orientation (extracted from ${}^{0}R_n$) can be converted to Euler
angles or quaternions using the formulas in Sections~2 and~6.
\end{keybox}

\begin{lstlisting}[style=matlab, caption={Forward kinematics of a 3-DOF
arm using DH parameters in MATLAB.}]
function T = DH(theta, d, a, alpha)
% Build 4x4 DH transformation matrix
  ct = cos(theta); st = sin(theta);
  ca = cos(alpha); sa = sin(alpha);
  T = [ct, -st*ca,  st*sa, a*ct;
       st,  ct*ca, -ct*sa, a*st;
        0,     sa,     ca,    d;
        0,      0,      0,    1];
end

% Example: 3-link planar arm (all revolute, alpha=0, d=0)
L = [0.5, 0.4, 0.3];          % link lengths [m]
q = deg2rad([30, 45, -20]);    % joint angles [rad]

T_total = eye(4);
for i = 1:3
    T_total = T_total * DH(q(i), 0, L(i), 0);
end

p_ee = T_total(1:3, 4);           % end effector position
R_ee = T_total(1:3, 1:3);         % end effector orientation
eul  = rotm2eul(R_ee, 'ZYX');     % convert to Euler angles
\end{lstlisting}

\end{document}
